\clearpage\pagestyle{otMain}\setlength{\otBlocEspacement}{2.5mm}\fontsize{11}{15}\selectfont\corrigesChapitre{3}{Calculer avec les décimaux}
\begin{blocLivre}{cor-3-1}\begin{corrige}{1}{Additions posées}\demarche{Aligner les virgules.}
$12{,}70+3{,}85=\res{16{,}55}$ ; $45{,}08+9{,}90=\res{54{,}98}$. Estimations : 17 et 55.
\end{corrige}\end{blocLivre}
\begin{blocLivre}{cor-3-2}\begin{corrige}{2}{Soustractions posées}\demarche{Compléter par des zéros avant de soustraire.}
$20{,}00-7{,}35=\res{12{,}65}$ ; $8{,}40-2{,}76=\res{5{,}64}$. Contrôle : $12{,}65+7{,}35=20$.
\end{corrige}\end{blocLivre}
\begin{blocLivre}{cor-3-3}\begin{corrige}{3}{Par 10, 100 ou 1 000}\demarche{Déplacer les chiffres de 1, 2 ou 3 rangs.}
$\res{456}$ ; $\res{3{,}7}$ ; $\res{2\,500}$ ; $\res{0{,}372}$ ; $\res{0{,}005}$.
\end{corrige}\end{blocLivre}
\begin{blocLivre}{cor-3-4}\begin{corrige}{4}{Par 0,1 ; 0,01 ; 0,001}\demarche{Multiplier par 0,1 revient à diviser par 10.}
$\res{4{,}5}$ ; $\res{0{,}062}$ ; $\res{0{,}3}$ ; $\res{0{,}07}$ : on divise par 10, 100 ou $1\,000$.
\end{corrige}\end{blocLivre}
\begin{blocLivre}{cor-3-5}\begin{corrige}{5}{Premiers produits}\demarche{Multiplier sans virgule, puis compter les décimales.}
$23\times14=322$, deux décimales : $\res{3{,}22}$. $5\times8=40$, deux décimales : $\res{0{,}4}$. $12\times3=36$, une décimale : $\res{3{,}6}$.
\end{corrige}\end{blocLivre}
\begin{blocLivre}{cor-3-6}\begin{corrige}{6}{Avec un rectangle}\demarche{Découper le rectangle pour calculer.}
Partie de 3 cm sur $2{,}5$ cm : $7{,}5\ \mathrm{cm}^2$. Partie de $0{,}2$ cm sur $2{,}5$ cm : $0{,}5\ \mathrm{cm}^2$. Total : $3{,}2\times2{,}5=\res{8}$.
\label{fig-cor-3-6}\figureGeo[corrige]{cal-6}
\end{corrige}\end{blocLivre}
\begin{blocLivre}{cor-3-7}\begin{corrige}{7}{Placer la virgule}\demarche{Compter les décimales des deux facteurs.}
$\res{33{,}81}$ ; $\res{33{,}81}$ ; $\res{0{,}033\,81}$ ; $\res{3{,}381}$ : deux, deux, cinq puis trois décimales.
\end{corrige}\end{blocLivre}
\begin{blocLivre}{cor-3-8}\begin{corrige}{8}{Ordre de grandeur}\demarche{Arrondir les facteurs pour estimer.}
$20\times5=100$ : $\res{100{,}98}$. $0{,}5\times60=30$ : $\res{29{,}89}$.
\end{corrige}\end{blocLivre}
\begin{blocLivre}{cor-3-9}\begin{corrige}{9}{Calculer astucieusement}\demarche{Regrouper les facteurs qui donnent 1 ou 10.}
$2{,}5\times4=10$, donc $\res{73}$. $0{,}5\times2=1$, donc $\res{36{,}8}$. $1{,}25\times8=10$, donc $\res{31}$.
\end{corrige}\end{blocLivre}
\begin{blocLivre}{cor-3-10}\begin{corrige}{10}{Plus petit ou plus grand ?}\demarche{Comparer le second facteur à 1.}
$0{,}9<1$ : plus petit, $\res{21{,}6}$. $1{,}2>1$ : plus grand, $\res{28{,}8}$. $0{,}05<1$ : plus petit, $\res{1{,}2}$.
\end{corrige}\end{blocLivre}
\begin{blocLivre}{cor-3-11}\begin{corrige}{11}{Les courses}\demarche{Calculer chaque achat, puis la monnaie.}
Pommes : $3\times2{,}45=7{,}35$ €. Raisin : $1{,}5\times4{,}80=7{,}20$ €. Total : $14{,}55$ €. Monnaie : $20-14{,}55=\res{5{,}45}$ €.
\end{corrige}\end{blocLivre}
\begin{blocLivre}{cor-3-12}\begin{corrige}{12}{Le terrain}\demarche{Estimer, puis poser le produit des entiers.}
Estimation : $25\times13=325\ \mathrm{m}^2$. Calcul : $245\times128=31\,360$, deux décimales : $\res{313{,}6\ \mathrm{m}^2}$.
\end{corrige}\end{blocLivre}
\begin{blocLivre}{cor-3-13}\begin{corrige}{13}{Le plein}\demarche{Compter les centaines de kilomètres.}
350 km, c’est $3{,}5$ fois 100 km : $3{,}5\times5{,}8=\res{20{,}3}$ L. Coût : $20{,}3\times1{,}80=\res{36{,}54}$ €.
\end{corrige}\end{blocLivre}
\begin{blocLivre}{cor-3-14}\begin{corrige}{14}{Le ruban}\demarche{Multiplier, puis soustraire.}
$12\times1{,}85=22{,}2$ m. Il reste $25-22{,}2=\res{2{,}8}$ m.
\end{corrige}\end{blocLivre}
\begin{blocLivre}{cor-3-15}\begin{corrige}{15}{L’erreur de virgule}\demarche{Contrôler avec le sens du calcul.}
Multiplier par $0{,}5$, c’est prendre la moitié : le résultat doit être plus petit que $3{,}6$. $36\times5=180$ et il faut deux décimales : $\res{1{,}8}$.
\end{corrige}\end{blocLivre}
